% Options for packages loaded elsewhere
% Options for packages loaded elsewhere
\PassOptionsToPackage{unicode}{hyperref}
\PassOptionsToPackage{hyphens}{url}
\PassOptionsToPackage{dvipsnames,svgnames,x11names}{xcolor}
%
\documentclass[
  brazilian,
  a4paper,
  DIV=11,
  numbers=noendperiod]{scrreprt}
\usepackage{xcolor}
\usepackage[top=30mm,left=30mm,bottom=25mm,right=25mm]{geometry}
\usepackage{amsmath,amssymb}
\setcounter{secnumdepth}{5}
\usepackage{iftex}
\ifPDFTeX
  \usepackage[T1]{fontenc}
  \usepackage[utf8]{inputenc}
  \usepackage{textcomp} % provide euro and other symbols
\else % if luatex or xetex
  \usepackage{unicode-math} % this also loads fontspec
  \defaultfontfeatures{Scale=MatchLowercase}
  \defaultfontfeatures[\rmfamily]{Ligatures=TeX,Scale=1}
\fi
\usepackage{lmodern}
\ifPDFTeX\else
  % xetex/luatex font selection
\fi
% Use upquote if available, for straight quotes in verbatim environments
\IfFileExists{upquote.sty}{\usepackage{upquote}}{}
\IfFileExists{microtype.sty}{% use microtype if available
  \usepackage[]{microtype}
  \UseMicrotypeSet[protrusion]{basicmath} % disable protrusion for tt fonts
}{}
\makeatletter
\@ifundefined{KOMAClassName}{% if non-KOMA class
  \IfFileExists{parskip.sty}{%
    \usepackage{parskip}
  }{% else
    \setlength{\parindent}{0pt}
    \setlength{\parskip}{6pt plus 2pt minus 1pt}}
}{% if KOMA class
  \KOMAoptions{parskip=half}}
\makeatother
% Make \paragraph and \subparagraph free-standing
\makeatletter
\ifx\paragraph\undefined\else
  \let\oldparagraph\paragraph
  \renewcommand{\paragraph}{
    \@ifstar
      \xxxParagraphStar
      \xxxParagraphNoStar
  }
  \newcommand{\xxxParagraphStar}[1]{\oldparagraph*{#1}\mbox{}}
  \newcommand{\xxxParagraphNoStar}[1]{\oldparagraph{#1}\mbox{}}
\fi
\ifx\subparagraph\undefined\else
  \let\oldsubparagraph\subparagraph
  \renewcommand{\subparagraph}{
    \@ifstar
      \xxxSubParagraphStar
      \xxxSubParagraphNoStar
  }
  \newcommand{\xxxSubParagraphStar}[1]{\oldsubparagraph*{#1}\mbox{}}
  \newcommand{\xxxSubParagraphNoStar}[1]{\oldsubparagraph{#1}\mbox{}}
\fi
\makeatother


\usepackage{longtable,booktabs,array}
\newcounter{none} % for unnumbered tables
\usepackage{calc} % for calculating minipage widths
% Correct order of tables after \paragraph or \subparagraph
\usepackage{etoolbox}
\makeatletter
\patchcmd\longtable{\par}{\if@noskipsec\mbox{}\fi\par}{}{}
\makeatother
% Allow footnotes in longtable head/foot
\IfFileExists{footnotehyper.sty}{\usepackage{footnotehyper}}{\usepackage{footnote}}
\makesavenoteenv{longtable}
\usepackage{graphicx}
\makeatletter
\newsavebox\pandoc@box
\newcommand*\pandocbounded[1]{% scales image to fit in text height/width
  \sbox\pandoc@box{#1}%
  \Gscale@div\@tempa{\textheight}{\dimexpr\ht\pandoc@box+\dp\pandoc@box\relax}%
  \Gscale@div\@tempb{\linewidth}{\wd\pandoc@box}%
  \ifdim\@tempb\p@<\@tempa\p@\let\@tempa\@tempb\fi% select the smaller of both
  \ifdim\@tempa\p@<\p@\scalebox{\@tempa}{\usebox\pandoc@box}%
  \else\usebox{\pandoc@box}%
  \fi%
}
% Set default figure placement to htbp
\def\fps@figure{htbp}
\makeatother


% definitions for citeproc citations
\NewDocumentCommand\citeproctext{}{}
\NewDocumentCommand\citeproc{mm}{%
  \begingroup\def\citeproctext{#2}\cite{#1}\endgroup}
\makeatletter
 % allow citations to break across lines
 \let\@cite@ofmt\@firstofone
 % avoid brackets around text for \cite:
 \def\@biblabel#1{}
 \def\@cite#1#2{{#1\if@tempswa , #2\fi}}
\makeatother
\newlength{\cslhangindent}
\setlength{\cslhangindent}{1.5em}
\newlength{\csllabelwidth}
\setlength{\csllabelwidth}{3em}
\newenvironment{CSLReferences}[2] % #1 hanging-indent, #2 entry-spacing
 {\begin{list}{}{%
  \setlength{\itemindent}{0pt}
  \setlength{\leftmargin}{0pt}
  \setlength{\parsep}{0pt}
  % turn on hanging indent if param 1 is 1
  \ifodd #1
   \setlength{\leftmargin}{\cslhangindent}
   \setlength{\itemindent}{-1\cslhangindent}
  \fi
  % set entry spacing
  \setlength{\itemsep}{#2\baselineskip}}}
 {\end{list}}
\usepackage{calc}
\newcommand{\CSLBlock}[1]{\hfill\break\parbox[t]{\linewidth}{\strut\ignorespaces#1\strut}}
\newcommand{\CSLLeftMargin}[1]{\parbox[t]{\csllabelwidth}{\strut#1\strut}}
\newcommand{\CSLRightInline}[1]{\parbox[t]{\linewidth - \csllabelwidth}{\strut#1\strut}}
\newcommand{\CSLIndent}[1]{\hspace{\cslhangindent}#1}

\ifLuaTeX
\usepackage[bidi=basic,shorthands=off]{babel}
\else
\usepackage[bidi=default,shorthands=off]{babel}
\fi
\ifLuaTeX
  \usepackage{selnolig} % disable illegal ligatures
\fi


\setlength{\emergencystretch}{3em} % prevent overfull lines

\providecommand{\tightlist}{%
  \setlength{\itemsep}{0pt}\setlength{\parskip}{0pt}}



 


\usepackage{booktabs}
\usepackage{longtable}
\usepackage{float}
\usepackage{amsmath}
\usepackage{amssymb}
\KOMAoption{captions}{tableheading}
\makeatletter
\@ifpackageloaded{bookmark}{}{\usepackage{bookmark}}
\makeatother
\makeatletter
\@ifpackageloaded{caption}{}{\usepackage{caption}}
\AtBeginDocument{%
\ifdefined\contentsname
  \renewcommand*\contentsname{Índice}
\else
  \newcommand\contentsname{Índice}
\fi
\ifdefined\listfigurename
  \renewcommand*\listfigurename{Lista de Figuras}
\else
  \newcommand\listfigurename{Lista de Figuras}
\fi
\ifdefined\listtablename
  \renewcommand*\listtablename{Lista de Tabelas}
\else
  \newcommand\listtablename{Lista de Tabelas}
\fi
\ifdefined\figurename
  \renewcommand*\figurename{Figura}
\else
  \newcommand\figurename{Figura}
\fi
\ifdefined\tablename
  \renewcommand*\tablename{Tabela}
\else
  \newcommand\tablename{Tabela}
\fi
}
\@ifpackageloaded{float}{}{\usepackage{float}}
\floatstyle{ruled}
\@ifundefined{c@chapter}{\newfloat{codelisting}{h}{lop}}{\newfloat{codelisting}{h}{lop}[chapter]}
\floatname{codelisting}{Listagem}
\newcommand*\listoflistings{\listof{codelisting}{Lista de Listagens}}
\makeatother
\makeatletter
\makeatother
\makeatletter
\@ifpackageloaded{caption}{}{\usepackage{caption}}
\@ifpackageloaded{subcaption}{}{\usepackage{subcaption}}
\makeatother
\usepackage{bookmark}
\IfFileExists{xurl.sty}{\usepackage{xurl}}{} % add URL line breaks if available
\urlstyle{same}
\hypersetup{
  pdftitle={Inteligência Artificial na Engenharia Aeroespacial},
  pdfauthor={Prof.~Dr.~Victor Rafael Rezende Celestino},
  pdflang={pt-BR},
  colorlinks=true,
  linkcolor={blue},
  filecolor={Maroon},
  citecolor={Blue},
  urlcolor={Blue},
  pdfcreator={LaTeX via pandoc}}


\title{Inteligência Artificial na Engenharia Aeroespacial}
\usepackage{etoolbox}
\makeatletter
\providecommand{\subtitle}[1]{% add subtitle to \maketitle
  \apptocmd{\@title}{\par {\large #1 \par}}{}{}
}
\makeatother
\subtitle{Minicurso panorâmico-aplicado --- SEMUNI 2026}
\author{Prof.~Dr.~Victor Rafael Rezende Celestino}
\date{20 de setembro de 2026}
\begin{document}
\maketitle

\renewcommand*\contentsname{Índice}
{
\setcounter{tocdepth}{2}
\tableofcontents
}

\bookmarksetup{startatroot}

\chapter*{Boas-vindas}\label{boas-vindas}
\addcontentsline{toc}{chapter}{Boas-vindas}

\markboth{Boas-vindas}{Boas-vindas}

Este material organiza o minicurso \textbf{Inteligência Artificial na
Engenharia Aeroespacial}, proposto para a \textbf{SEMUNI 2026 ---
FCTE/UnB}.

O minicurso tem caráter \textbf{panorâmico-aplicado}: em quatro horas,
seu propósito não é formar especialistas em cada arquitetura, mas
apresentar uma visão integrada de como inteligência artificial,
simulação, modelos substitutos, grafos e otimização podem apoiar
problemas de engenharia aeroespacial.

\section*{Informações gerais}\label{informauxe7uxf5es-gerais}
\addcontentsline{toc}{section}{Informações gerais}

\markright{Informações gerais}

{\def\LTcaptype{none} % do not increment counter
\begin{longtable}[]{@{}
  >{\raggedright\arraybackslash}p{(\linewidth - 2\tabcolsep) * \real{0.5000}}
  >{\raggedright\arraybackslash}p{(\linewidth - 2\tabcolsep) * \real{0.5000}}@{}}
\toprule\noalign{}
\begin{minipage}[b]{\linewidth}\raggedright
Item
\end{minipage} & \begin{minipage}[b]{\linewidth}\raggedright
Descrição
\end{minipage} \\
\midrule\noalign{}
\endhead
\bottomrule\noalign{}
\endlastfoot
Evento & SEMUNI 2026 --- Semana Universitária da FCTE/UnB \\
Período previsto & 20 a 25 de setembro de 2026 \\
Título & Inteligência Artificial na Engenharia Aeroespacial \\
Carga horária & 4 horas \\
Professor responsável & Prof.~Dr.~Victor Rafael Rezende Celestino \\
Formato & Exposição dialogada, roteiro navegável, demonstrações
computacionais e discussão de aplicações \\
Natureza & Panorama aplicado sobre IA, simulação, otimização e projeto
aeroespacial \\
\end{longtable}
}

\section*{Síntese executiva}\label{suxedntese-executiva}
\addcontentsline{toc}{section}{Síntese executiva}

\markright{Síntese executiva}

A engenharia aeroespacial vem incorporando técnicas de aprendizado
profundo, modelos generativos, redes neurais de grafo, simulação
orientada por dados e otimização baseada em modelos. Essas técnicas
ampliam a capacidade de modelar fenômenos, acelerar simulações, explorar
alternativas de projeto e apoiar decisões em sistemas complexos, sem
substituir os fundamentos físicos, matemáticos e experimentais da
engenharia.

O minicurso articula quatro eixos:

\begin{enumerate}
\def\labelenumi{\arabic{enumi}.}
\tightlist
\item
  \textbf{Deep Learning, MLPs e PINNs};
\item
  \textbf{IA generativa, dados sintéticos e modelos substitutos};
\item
  \textbf{GNNs, otimização combinatória, MINLP e MPC};
\item
  \textbf{IA em otimização multidisciplinar e AeroDesign}.
\end{enumerate}

\section*{Público-alvo}\label{puxfablico-alvo}
\addcontentsline{toc}{section}{Público-alvo}

\markright{Público-alvo}

Estudantes de graduação e pós-graduação em Engenharia Aeroespacial,
Engenharia Automotiva, Engenharia de Energia, Engenharia Eletrônica,
Engenharia de Software, Computação, Ciência de Dados, Administração,
Pesquisa Operacional e áreas afins; professores, pesquisadores e
técnicos interessados na interseção entre inteligência artificial e
aplicações aeroespaciais.

\section*{Pré-requisitos
desejáveis}\label{pruxe9-requisitos-desejuxe1veis}
\addcontentsline{toc}{section}{Pré-requisitos desejáveis}

\markright{Pré-requisitos desejáveis}

\begin{itemize}
\tightlist
\item
  Noções básicas de cálculo, álgebra linear e física;
\item
  conhecimento introdutório de Python;
\item
  fundamentos de engenharia;
\item
  familiaridade inicial com aprendizado de máquina, embora não
  obrigatória.
\end{itemize}

\section*{Cronograma de 4 horas}\label{cronograma-de-4-horas}
\addcontentsline{toc}{section}{Cronograma de 4 horas}

\markright{Cronograma de 4 horas}

{\def\LTcaptype{none} % do not increment counter
\begin{longtable}[]{@{}
  >{\raggedleft\arraybackslash}p{(\linewidth - 4\tabcolsep) * \real{0.4000}}
  >{\raggedright\arraybackslash}p{(\linewidth - 4\tabcolsep) * \real{0.3000}}
  >{\raggedright\arraybackslash}p{(\linewidth - 4\tabcolsep) * \real{0.3000}}@{}}
\toprule\noalign{}
\begin{minipage}[b]{\linewidth}\raggedleft
Tempo
\end{minipage} & \begin{minipage}[b]{\linewidth}\raggedright
Bloco
\end{minipage} & \begin{minipage}[b]{\linewidth}\raggedright
Conteúdo principal
\end{minipage} \\
\midrule\noalign{}
\endhead
\bottomrule\noalign{}
\endlastfoot
0:00--0:15 & Abertura & Motivação, escopo panorâmico-aplicado e mapa do
repositório \\
0:15--1:00 & Módulo 1 & MLPs, deep learning, PINNs e demonstração em
PyTorch/DeepXDE \\
1:00--2:00 & Módulo 2 & IA generativa, dados sintéticos, geometrias e
surrogate models \\
2:00--3:00 & Módulo 3 & GNNs, grafos bipartidos, MINLP, MPC e PyTorch
Geometric/DGL \\
3:00--3:50 & Módulo 4 & IA, MDO, OpenMDAO, OpenVSP e estudo inspirado no
AeroDesign \\
3:50--4:00 & Fechamento & Síntese, bibliografia, repositório e
oportunidades de projetos \\
\end{longtable}
}

\section*{Regra operacional do curso}\label{regra-operacional-do-curso}
\addcontentsline{toc}{section}{Regra operacional do curso}

\markright{Regra operacional do curso}

Cada módulo terá \textbf{uma demonstração principal}. Conteúdos
adicionais serão apresentados como leituras, desafios e extensões para
estudo posterior.

\section*{Como navegar}\label{como-navegar}
\addcontentsline{toc}{section}{Como navegar}

\markright{Como navegar}

Use o menu lateral para acessar os quatro módulos. Em cada módulo,
procure as seções:

\begin{itemize}
\tightlist
\item
  problema de engenharia;
\item
  classe de modelo de IA;
\item
  fundamentos mínimos;
\item
  roteiro de demonstração ao vivo;
\item
  links para abrir durante a aula;
\item
  notebook associado;
\item
  cuidados de engenharia;
\item
  projetos derivados.
\end{itemize}

\bookmarksetup{startatroot}

\chapter{Apresentação e filosofia do
curso}\label{apresentauxe7uxe3o-e-filosofia-do-curso}

\section{Tese central}\label{tese-central}

A inteligência artificial deve ser compreendida como \textbf{acelerador,
aproximador e organizador de conhecimento} em engenharia aeroespacial,
não como substituta dos princípios físicos, matemáticos e experimentais
que sustentam a engenharia.

Em problemas aeroespaciais, modelos de IA podem:

\begin{itemize}
\tightlist
\item
  aproximar funções de alto custo computacional;
\item
  apoiar a exploração de espaços de projeto;
\item
  reduzir o tempo de simulações preliminares;
\item
  extrair padrões de dados experimentais, numéricos ou operacionais;
\item
  representar sistemas complexos por grafos;
\item
  apoiar decisões em otimização e controle.
\end{itemize}

Contudo, esses modelos devem ser avaliados com rigor quanto a:

\begin{itemize}
\tightlist
\item
  domínio de validade;
\item
  extrapolação;
\item
  incerteza;
\item
  consistência física;
\item
  interpretabilidade;
\item
  robustez computacional;
\item
  rastreabilidade dos dados e hipóteses.
\end{itemize}

\section{Diretriz pedagógica dos
módulos}\label{diretriz-pedaguxf3gica-dos-muxf3dulos}

Cada módulo será guiado por três perguntas:

\begin{enumerate}
\def\labelenumi{\arabic{enumi}.}
\tightlist
\item
  \textbf{Qual problema de engenharia está sendo tratado?}
\item
  \textbf{Qual classe de modelo de IA é adequada e por quê?}
\item
  \textbf{Quais cuidados de validação, extrapolação e consistência
  física devem ser observados?}
\end{enumerate}

Essa estrutura busca conectar intuição física, formulação matemática e
demonstração computacional.

\section{Escopo e limites}\label{escopo-e-limites}

Este minicurso é uma introdução aplicada. Em quatro horas, não é
realista treinar profundamente os participantes em PyTorch, DeepXDE,
PyTorch Geometric, DGL, OpenMDAO, OpenAeroStruct e OpenVSP. O objetivo é
oferecer um mapa conceitual, um conjunto de exemplos e um repositório
público para aprofundamento.

\part{Módulos do curso}

\chapter{Módulo 1 --- Deep Learning, MLPs e
PINNs}\label{muxf3dulo-1-deep-learning-mlps-e-pinns}

\section{1. Objetivos de aprendizagem}\label{objetivos-de-aprendizagem}

Ao final deste módulo, espera-se que o participante seja capaz de:

\begin{itemize}
\tightlist
\item
  reconhecer o problema de engenharia associado ao módulo;
\item
  compreender a intuição da classe de modelo de IA apresentada;
\item
  identificar uma demonstração computacional mínima;
\item
  avaliar riscos de extrapolação, validação e consistência física;
\item
  localizar recursos públicos para aprofundamento.
\end{itemize}

\section{2. Problema de engenharia}\label{problema-de-engenharia}

Predição e aproximação de respostas estruturais ou aerodinâmicas em
problemas simplificados, especialmente quando há necessidade de combinar
dados com restrições físicas.

\section{3. Classe de modelo de IA}\label{classe-de-modelo-de-ia}

Perceptrons multicamadas, redes neurais profundas e redes neurais
informadas pela física. O foco é compreender a progressão de MLPs para
arquiteturas modernas e como PINNs incorporam resíduos de equações
diferenciais na função de perda.

\section{4. Fundamentos conceituais}\label{fundamentos-conceituais}

MLP, função de ativação, treinamento por retropropagação, autograd,
regularização, overfitting, modelos fundacionais e aprendizado
físico-informado.

\section{5. Formulação matemática
mínima}\label{formulauxe7uxe3o-matemuxe1tica-muxednima}

Uma PINN pode ser vista como uma rede neural (\(\hat{u}_\theta(x,t)\))
treinada por uma função de perda composta:

\[
\mathcal{L}(\theta) =
\lambda_d \mathcal{L}_{dados} +
\lambda_f \mathcal{L}_{fisica} +
\lambda_b \mathcal{L}_{contorno} +
\lambda_i \mathcal{L}_{inicial}.
\]

A contribuição física geralmente envolve o resíduo de uma equação
diferencial:

\[
r_\theta(x,t) = \mathcal{N}[\hat{u}_\theta](x,t),
\]

onde (\(\mathcal{N}\)) representa o operador diferencial do problema.

\section{6. Mapa de demonstração ao
vivo}\label{mapa-de-demonstrauxe7uxe3o-ao-vivo}

{\def\LTcaptype{none} % do not increment counter
\begin{longtable}[]{@{}
  >{\raggedleft\arraybackslash}p{(\linewidth - 4\tabcolsep) * \real{0.4000}}
  >{\raggedright\arraybackslash}p{(\linewidth - 4\tabcolsep) * \real{0.3000}}
  >{\raggedright\arraybackslash}p{(\linewidth - 4\tabcolsep) * \real{0.3000}}@{}}
\toprule\noalign{}
\begin{minipage}[b]{\linewidth}\raggedleft
Passo
\end{minipage} & \begin{minipage}[b]{\linewidth}\raggedright
Ação em sala
\end{minipage} & \begin{minipage}[b]{\linewidth}\raggedright
Mensagem didática
\end{minipage} \\
\midrule\noalign{}
\endhead
\bottomrule\noalign{}
\endlastfoot
1 & Abrir a página do módulo no GitHub Pages & Apresentar o problema de
engenharia \\
2 & Abrir o notebook associado & Mostrar o fluxo computacional mínimo \\
3 & Executar ou comentar as células principais & Relacionar código,
modelo e hipótese física \\
4 & Alterar um parâmetro simples & Evidenciar sensibilidade e limites do
modelo \\
5 & Discutir validação & Evitar uso acrítico da IA em engenharia \\
\end{longtable}
}

\section{7. Links para abrir durante a
aula}\label{links-para-abrir-durante-a-aula}

\begin{itemize}
\tightlist
\item
  PyTorch Documentation;
\item
  DeepXDE Documentation;
\item
  NASA Airfoil Self-Noise Dataset;
\item
  repositórios introdutórios de PINNs em PyTorch.
\end{itemize}

\section{8. Notebook associado}\label{notebook-associado}

\begin{quote}
Notebook principal: \texttt{notebooks/01\_mlp\_pinn\_pytorch.ipynb}
\end{quote}

Quando o repositório estiver publicado no GitHub, este espaço receberá o
botão direto para abertura no Google Colab.

\section{9. Cuidados de engenharia}\label{cuidados-de-engenharia}

\begin{itemize}
\tightlist
\item
  Não interpretar baixa perda de treinamento como validade física geral;
\item
  verificar condições de contorno;
\item
  analisar extrapolação fora do domínio treinado;
\item
  comparar com solução analítica, numérica ou experimental quando
  disponível.
\end{itemize}

\section{10. Questões de discussão}\label{questuxf5es-de-discussuxe3o}

\begin{itemize}
\tightlist
\item
  Quando uma PINN é preferível a uma rede puramente orientada por dados?
\item
  Que riscos surgem quando a física é mal especificada?
\item
  Como escolher pesos entre dados, resíduo físico e condições de
  contorno?
\end{itemize}

\section{11. Possíveis projetos
derivados}\label{possuxedveis-projetos-derivados}

\begin{itemize}
\tightlist
\item
  PINN para viga Euler-Bernoulli simplificada;
\item
  MLP para aproximação de coeficientes aerodinâmicos;
\item
  comparação entre regressão clássica, MLP e PINN.
\end{itemize}

\section{12. Leituras recomendadas}\label{leituras-recomendadas}

\begin{itemize}
\tightlist
\item
  Goodfellow et al. (2016)
\item
  LeCun et al. (2015)
\item
  Raissi et al. (2019)
\item
  Karniadakis et al. (2021)
\item
  Brunton e Kutz (2022)
\end{itemize}

\chapter{Módulo 2 --- IA generativa, dados sintéticos e modelos
substitutos}\label{muxf3dulo-2-ia-generativa-dados-sintuxe9ticos-e-modelos-substitutos}

\section{1. Objetivos de
aprendizagem}\label{objetivos-de-aprendizagem-1}

Ao final deste módulo, espera-se que o participante seja capaz de:

\begin{itemize}
\tightlist
\item
  reconhecer o problema de engenharia associado ao módulo;
\item
  compreender a intuição da classe de modelo de IA apresentada;
\item
  identificar uma demonstração computacional mínima;
\item
  avaliar riscos de extrapolação, validação e consistência física;
\item
  localizar recursos públicos para aprofundamento.
\end{itemize}

\section{2. Problema de engenharia}\label{problema-de-engenharia-1}

Geração de dados sintéticos, exploração de geometrias e construção de
modelos substitutos para reduzir custo computacional em simulações
aeroespaciais.

\section{3. Classe de modelo de IA}\label{classe-de-modelo-de-ia-1}

Modelos generativos, incluindo VAEs, GANs, modelos de difusão, modelos
autorregressivos e mecanismos de atenção. O módulo enfatiza também
surrogate models como aproximadores de simulações caras.

\section{4. Fundamentos conceituais}\label{fundamentos-conceituais-1}

Distribuição de dados, espaço latente, codificador, decodificador,
discriminador, difusão, atenção, plausibilidade física, aumento de dados
e aprendizado ativo.

\section{5. Formulação matemática
mínima}\label{formulauxe7uxe3o-matemuxe1tica-muxednima-1}

Um modelo substituto busca aproximar uma função de alto custo:

\[
y = f(x) \quad \Rightarrow \quad \hat{y} = \hat{f}_\theta(x),
\]

em que (\(f\)) pode representar uma simulação CFD, análise estrutural ou
cálculo de desempenho.

Em um VAE, o treinamento combina reconstrução e regularização
probabilística:

\[
\mathcal{L} =
\mathbb{E}_{q_\phi(z|x)}[\log p_\theta(x|z)] -
D_{KL}(q_\phi(z|x)\|p(z)).
\]

\section{6. Mapa de demonstração ao
vivo}\label{mapa-de-demonstrauxe7uxe3o-ao-vivo-1}

{\def\LTcaptype{none} % do not increment counter
\begin{longtable}[]{@{}
  >{\raggedleft\arraybackslash}p{(\linewidth - 4\tabcolsep) * \real{0.4000}}
  >{\raggedright\arraybackslash}p{(\linewidth - 4\tabcolsep) * \real{0.3000}}
  >{\raggedright\arraybackslash}p{(\linewidth - 4\tabcolsep) * \real{0.3000}}@{}}
\toprule\noalign{}
\begin{minipage}[b]{\linewidth}\raggedleft
Passo
\end{minipage} & \begin{minipage}[b]{\linewidth}\raggedright
Ação em sala
\end{minipage} & \begin{minipage}[b]{\linewidth}\raggedright
Mensagem didática
\end{minipage} \\
\midrule\noalign{}
\endhead
\bottomrule\noalign{}
\endlastfoot
1 & Abrir a página do módulo no GitHub Pages & Apresentar o problema de
engenharia \\
2 & Abrir o notebook associado & Mostrar o fluxo computacional mínimo \\
3 & Executar ou comentar as células principais & Relacionar código,
modelo e hipótese física \\
4 & Alterar um parâmetro simples & Evidenciar sensibilidade e limites do
modelo \\
5 & Discutir validação & Evitar uso acrítico da IA em engenharia \\
\end{longtable}
}

\section{7. Links para abrir durante a
aula}\label{links-para-abrir-durante-a-aula-1}

\begin{itemize}
\tightlist
\item
  NASA Airfoil Self-Noise Dataset;
\item
  UIUC Airfoil Data Site;
\item
  ShapeNet / ShapeNet Part;
\item
  PyTorch Documentation;
\item
  exemplos de surrogate modeling em Scikit-learn e PyTorch.
\end{itemize}

\section{8. Notebook associado}\label{notebook-associado-1}

\begin{quote}
Notebook principal:
\texttt{notebooks/02\_surrogate\_dados\_sinteticos.ipynb}
\end{quote}

Quando o repositório estiver publicado no GitHub, este espaço receberá o
botão direto para abertura no Google Colab.

\section{9. Cuidados de engenharia}\label{cuidados-de-engenharia-1}

\begin{itemize}
\tightlist
\item
  Dados sintéticos não substituem validação experimental;
\item
  geometrias geradas precisam respeitar restrições físicas e
  operacionais;
\item
  surrogate models podem falhar severamente fora do domínio de
  treinamento;
\item
  métricas estatísticas devem ser combinadas com julgamento de
  engenharia.
\end{itemize}

\section{10. Questões de discussão}\label{questuxf5es-de-discussuxe3o-1}

\begin{itemize}
\tightlist
\item
  Como validar dados sintéticos em engenharia?
\item
  Quando um surrogate model é suficiente para decisão preliminar?
\item
  Que tipo de erro é aceitável em projeto conceitual?
\end{itemize}

\section{11. Possíveis projetos
derivados}\label{possuxedveis-projetos-derivados-1}

\begin{itemize}
\tightlist
\item
  modelo substituto para ruído de aerofólio;
\item
  geração de amostras sintéticas de configurações de asa;
\item
  comparação entre Gaussian Process, Random Forest e MLP para regressão
  aerodinâmica.
\end{itemize}

\section{12. Leituras recomendadas}\label{leituras-recomendadas-1}

\begin{itemize}
\tightlist
\item
  Goodfellow et al. (2014)
\item
  Kingma e Welling (2013)
\item
  Ho et al. (2020)
\item
  Rombach et al. (2022)
\item
  Forrester et al. (2008)
\end{itemize}

\chapter{Módulo 3 --- GNNs, otimização, MINLP e
MPC}\label{muxf3dulo-3-gnns-otimizauxe7uxe3o-minlp-e-mpc}

\section{1. Objetivos de
aprendizagem}\label{objetivos-de-aprendizagem-2}

Ao final deste módulo, espera-se que o participante seja capaz de:

\begin{itemize}
\tightlist
\item
  reconhecer o problema de engenharia associado ao módulo;
\item
  compreender a intuição da classe de modelo de IA apresentada;
\item
  identificar uma demonstração computacional mínima;
\item
  avaliar riscos de extrapolação, validação e consistência física;
\item
  localizar recursos públicos para aprofundamento.
\end{itemize}

\section{2. Problema de engenharia}\label{problema-de-engenharia-2}

Representação de sistemas complexos, malhas, componentes, restrições e
problemas de otimização por meio de grafos, com aplicações em otimização
combinatória, MINLP e controle preditivo.

\section{3. Classe de modelo de IA}\label{classe-de-modelo-de-ia-2}

Redes neurais de grafo, grafos bipartidos, embeddings e message passing.
O módulo apresenta como problemas de otimização podem ser representados
por variáveis, restrições e relações estruturais.

\section{4. Fundamentos conceituais}\label{fundamentos-conceituais-2}

Nós, arestas, atributos, grafos bipartidos, embeddings, agregação de
vizinhança, message passing, warm start, predição de variáveis ativas e
suporte a decisões em horizonte móvel.

\section{5. Formulação matemática
mínima}\label{formulauxe7uxe3o-matemuxe1tica-muxednima-2}

Em uma GNN, uma etapa genérica de message passing pode ser escrita como:

\[
m_i^{(k)} = \sum_{j \in \mathcal{N}(i)}
\phi^{(k)}(h_i^{(k)}, h_j^{(k)}, e_{ij}),
\]

\[
h_i^{(k+1)} =
\psi^{(k)}(h_i^{(k)}, m_i^{(k)}),
\]

onde (\(h_i\)) é a representação do nó, (\(e_{ij}\)) representa
atributos da aresta e (\(\mathcal{N}(i)\)) é a vizinhança de (\(i\)).

\section{6. Mapa de demonstração ao
vivo}\label{mapa-de-demonstrauxe7uxe3o-ao-vivo-2}

{\def\LTcaptype{none} % do not increment counter
\begin{longtable}[]{@{}
  >{\raggedleft\arraybackslash}p{(\linewidth - 4\tabcolsep) * \real{0.4000}}
  >{\raggedright\arraybackslash}p{(\linewidth - 4\tabcolsep) * \real{0.3000}}
  >{\raggedright\arraybackslash}p{(\linewidth - 4\tabcolsep) * \real{0.3000}}@{}}
\toprule\noalign{}
\begin{minipage}[b]{\linewidth}\raggedleft
Passo
\end{minipage} & \begin{minipage}[b]{\linewidth}\raggedright
Ação em sala
\end{minipage} & \begin{minipage}[b]{\linewidth}\raggedright
Mensagem didática
\end{minipage} \\
\midrule\noalign{}
\endhead
\bottomrule\noalign{}
\endlastfoot
1 & Abrir a página do módulo no GitHub Pages & Apresentar o problema de
engenharia \\
2 & Abrir o notebook associado & Mostrar o fluxo computacional mínimo \\
3 & Executar ou comentar as células principais & Relacionar código,
modelo e hipótese física \\
4 & Alterar um parâmetro simples & Evidenciar sensibilidade e limites do
modelo \\
5 & Discutir validação & Evitar uso acrítico da IA em engenharia \\
\end{longtable}
}

\section{7. Links para abrir durante a
aula}\label{links-para-abrir-durante-a-aula-2}

\begin{itemize}
\tightlist
\item
  PyTorch Geometric Documentation;
\item
  DGL Documentation;
\item
  MIPLIB;
\item
  MILPBench;
\item
  benchmarks de facility location;
\item
  materiais introdutórios sobre aprendizado para otimização.
\end{itemize}

\section{8. Notebook associado}\label{notebook-associado-2}

\begin{quote}
Notebook principal: \texttt{notebooks/03\_gnn\_bipartite\_minlp.ipynb}
\end{quote}

Quando o repositório estiver publicado no GitHub, este espaço receberá o
botão direto para abertura no Google Colab.

\section{9. Cuidados de engenharia}\label{cuidados-de-engenharia-2}

\begin{itemize}
\tightlist
\item
  GNNs podem capturar padrões estruturais, mas não garantem otimalidade;
\item
  decisões geradas por IA devem ser verificadas por solucionadores ou
  regras de viabilidade;
\item
  benchmarks de otimização nem sempre representam sistemas aeroespaciais
  reais;
\item
  MPC exige tratamento rigoroso de estabilidade, restrições e robustez.
\end{itemize}

\section{10. Questões de discussão}\label{questuxf5es-de-discussuxe3o-2}

\begin{itemize}
\tightlist
\item
  Como representar um problema de otimização como grafo bipartido?
\item
  GNN deve substituir ou auxiliar o solver?
\item
  Como validar uma política aproximada em MPC?
\end{itemize}

\section{11. Possíveis projetos
derivados}\label{possuxedveis-projetos-derivados-2}

\begin{itemize}
\tightlist
\item
  grafo bipartido para pequeno problema MILP;
\item
  classificação de variáveis candidatas a zero/um;
\item
  warm start orientado por aprendizado;
\item
  GNN para rede simplificada de componentes aeroespaciais.
\end{itemize}

\section{12. Leituras recomendadas}\label{leituras-recomendadas-2}

\begin{itemize}
\tightlist
\item
  Kipf e Welling (2017)
\item
  Gilmer et al. (2017)
\item
  Battaglia et al. (2018)
\item
  Fey e Lenssen (2019)
\item
  Zhou et al. (2020)
\item
  Rawlings et al. (2017)
\end{itemize}

\chapter{Módulo 4 --- IA em MDO e projeto preliminar inspirado em
AeroDesign}\label{muxf3dulo-4-ia-em-mdo-e-projeto-preliminar-inspirado-em-aerodesign}

\section{1. Objetivos de
aprendizagem}\label{objetivos-de-aprendizagem-3}

Ao final deste módulo, espera-se que o participante seja capaz de:

\begin{itemize}
\tightlist
\item
  reconhecer o problema de engenharia associado ao módulo;
\item
  compreender a intuição da classe de modelo de IA apresentada;
\item
  identificar uma demonstração computacional mínima;
\item
  avaliar riscos de extrapolação, validação e consistência física;
\item
  localizar recursos públicos para aprofundamento.
\end{itemize}

\section{2. Problema de engenharia}\label{problema-de-engenharia-3}

Projeto preliminar de aeronaves de pequeno porte com disciplinas
acopladas: aerodinâmica, estruturas, massa, propulsão, desempenho,
estabilidade e restrições operacionais.

\section{3. Classe de modelo de IA}\label{classe-de-modelo-de-ia-3}

Modelos substitutos, otimização bayesiana, aprendizado ativo e
integração com frameworks de otimização multidisciplinar, especialmente
OpenMDAO e ferramentas de geometria conceitual.

\section{4. Fundamentos conceituais}\label{fundamentos-conceituais-3}

Variáveis de projeto, funções objetivo, restrições, disciplinas
acopladas, análise de sensibilidade, metamodelos, trade-offs e busca em
espaços de configuração.

\section{5. Formulação matemática
mínima}\label{formulauxe7uxe3o-matemuxe1tica-muxednima-3}

Um problema MDO simplificado pode ser escrito como:

\[
\min_{x} \; F(x)
\]

sujeito a:

\[
g_i(x) \leq 0,\quad i=1,\ldots,m,
\]

\[
h_j(x) = 0,\quad j=1,\ldots,p,
\]

em que (\(x\)) representa variáveis de projeto, (\(F(x)\)) é a função
objetivo e (\(g_i\), \(h_j\)) representam restrições de engenharia.

\section{6. Mapa de demonstração ao
vivo}\label{mapa-de-demonstrauxe7uxe3o-ao-vivo-3}

{\def\LTcaptype{none} % do not increment counter
\begin{longtable}[]{@{}
  >{\raggedleft\arraybackslash}p{(\linewidth - 4\tabcolsep) * \real{0.4000}}
  >{\raggedright\arraybackslash}p{(\linewidth - 4\tabcolsep) * \real{0.3000}}
  >{\raggedright\arraybackslash}p{(\linewidth - 4\tabcolsep) * \real{0.3000}}@{}}
\toprule\noalign{}
\begin{minipage}[b]{\linewidth}\raggedleft
Passo
\end{minipage} & \begin{minipage}[b]{\linewidth}\raggedright
Ação em sala
\end{minipage} & \begin{minipage}[b]{\linewidth}\raggedright
Mensagem didática
\end{minipage} \\
\midrule\noalign{}
\endhead
\bottomrule\noalign{}
\endlastfoot
1 & Abrir a página do módulo no GitHub Pages & Apresentar o problema de
engenharia \\
2 & Abrir o notebook associado & Mostrar o fluxo computacional mínimo \\
3 & Executar ou comentar as células principais & Relacionar código,
modelo e hipótese física \\
4 & Alterar um parâmetro simples & Evidenciar sensibilidade e limites do
modelo \\
5 & Discutir validação & Evitar uso acrítico da IA em engenharia \\
\end{longtable}
}

\section{7. Links para abrir durante a
aula}\label{links-para-abrir-durante-a-aula-3}

\begin{itemize}
\tightlist
\item
  OpenMDAO Documentation;
\item
  OpenAeroStruct Documentation;
\item
  OpenVSP Documentation;
\item
  UIUC Airfoil Data Site;
\item
  materiais públicos de AeroDesign e projeto conceitual de aeronaves.
\end{itemize}

\section{8. Notebook associado}\label{notebook-associado-3}

\begin{quote}
Notebook principal:
\texttt{notebooks/04\_mdo\_aerodesign\_openmdao.ipynb}
\end{quote}

Quando o repositório estiver publicado no GitHub, este espaço receberá o
botão direto para abertura no Google Colab.

\section{9. Cuidados de engenharia}\label{cuidados-de-engenharia-3}

\begin{itemize}
\tightlist
\item
  Modelos preliminares devem ser tratados como apoio à decisão, não como
  certificação;
\item
  acoplamentos entre disciplinas podem gerar efeitos não intuitivos;
\item
  restrições estruturais, operacionais e de segurança devem prevalecer
  sobre ganho de função objetivo;
\item
  otimização sem validação pode produzir soluções inviáveis.
\end{itemize}

\section{10. Questões de discussão}\label{questuxf5es-de-discussuxe3o-3}

\begin{itemize}
\tightlist
\item
  Como definir variáveis de projeto relevantes para uma aeronave de
  pequeno porte?
\item
  Qual é o papel de um surrogate model em MDO?
\item
  Como lidar com trade-offs entre carga paga, decolagem, massa e
  desempenho aerodinâmico?
\end{itemize}

\section{11. Possíveis projetos
derivados}\label{possuxedveis-projetos-derivados-3}

\begin{itemize}
\tightlist
\item
  MDO simplificado de asa para AeroDesign;
\item
  surrogate model para estimar carga paga;
\item
  estudo de sensibilidade de alongamento e área de asa;
\item
  integração OpenVSP--OpenMDAO em projeto conceitual.
\end{itemize}

\section{12. Leituras recomendadas}\label{leituras-recomendadas-3}

\begin{itemize}
\tightlist
\item
  Martins e Ning (2021)
\item
  Gray et al. (2019)
\item
  Jasa et al. (2018)
\item
  Forrester et al. (2008)
\item
  Brunton e Kutz (2022)
\end{itemize}

\part{Recursos e fechamento}

\chapter{Recursos computacionais, datasets e
repositórios}\label{recursos-computacionais-datasets-e-reposituxf3rios}

\section{Bibliotecas e frameworks}\label{bibliotecas-e-frameworks}

{\def\LTcaptype{none} % do not increment counter
\begin{longtable}[]{@{}
  >{\raggedright\arraybackslash}p{(\linewidth - 2\tabcolsep) * \real{0.5000}}
  >{\raggedright\arraybackslash}p{(\linewidth - 2\tabcolsep) * \real{0.5000}}@{}}
\toprule\noalign{}
\begin{minipage}[b]{\linewidth}\raggedright
Recurso
\end{minipage} & \begin{minipage}[b]{\linewidth}\raggedright
Uso no minicurso
\end{minipage} \\
\midrule\noalign{}
\endhead
\bottomrule\noalign{}
\endlastfoot
Python, NumPy, SciPy, Pandas, Matplotlib, Scikit-learn & Manipulação de
dados, modelagem, regressão e visualização \\
PyTorch & Construção de MLPs, redes profundas e protótipos de PINNs \\
DeepXDE & Alternativa especializada para redes neurais informadas pela
física \\
PyTorch Geometric & Construção e treinamento de GNNs \\
DGL & Alternativa para modelagem e treinamento de GNNs \\
OpenMDAO & Análise e otimização multidisciplinar \\
OpenAeroStruct & Otimização aeroestrutural acoplada \\
OpenVSP & Modelagem paramétrica e análise conceitual de aeronaves \\
Pyomo e SciPy Optimize & Modelagem e solução de problemas de
otimização \\
\end{longtable}
}

\section{Datasets e repositórios
públicos}\label{datasets-e-reposituxf3rios-puxfablicos}

{\def\LTcaptype{none} % do not increment counter
\begin{longtable}[]{@{}
  >{\raggedright\arraybackslash}p{(\linewidth - 2\tabcolsep) * \real{0.5000}}
  >{\raggedright\arraybackslash}p{(\linewidth - 2\tabcolsep) * \real{0.5000}}@{}}
\toprule\noalign{}
\begin{minipage}[b]{\linewidth}\raggedright
Recurso
\end{minipage} & \begin{minipage}[b]{\linewidth}\raggedright
Uso didático
\end{minipage} \\
\midrule\noalign{}
\endhead
\bottomrule\noalign{}
\endlastfoot
NASA Airfoil Self-Noise & Regressão aplicada à aerodinâmica/acústica de
aerofólios \\
UIUC Airfoil Data Site & Coordenadas de aerofólios para estudos
preliminares de geometria \\
NPARC Alliance / NASA CFD resources & Discussão de validação
computacional e fluidodinâmica \\
ShapeNet / ShapeNet Part & Representação geométrica, nuvens de pontos e
IA generativa \\
MIPLIB & Benchmarks de programação inteira mista \\
MILPBench e facility location benchmarks & Aprendizado para otimização e
geração de warm starts \\
Repositórios de PINNs em PyTorch e DeepXDE & Exemplos de EDOs, EDPs e
sistemas físicos \\
Repositórios OpenMDAO e OpenAeroStruct & Exemplos de MDO, análise
aeroestrutural e otimização \\
\end{longtable}
}

\section{Política de dados}\label{poluxedtica-de-dados}

Este repositório não redistribui datasets de terceiros por padrão.
Sempre que possível, serão fornecidos links, instruções de acesso e
scripts de download, respeitando as licenças originais.

\chapter{Referências}\label{referuxeancias}

\protect\phantomsection\label{refs}
\begin{CSLReferences}{1}{1}
\bibitem[\citeproctext]{ref-battaglia2018relational}
Battaglia, Peter W., Jessica B. Hamrick, Victor Bapst, et al. 2018.
\emph{Relational inductive biases, deep learning, and graph networks}.
\url{https://arxiv.org/abs/1806.01261}.

\bibitem[\citeproctext]{ref-brunton2022data}
Brunton, Steven L., e J. Nathan Kutz. 2022. \emph{Data-Driven Science
and Engineering: Machine Learning, Dynamical Systems, and Control}. 2º
ed. Cambridge University Press.

\bibitem[\citeproctext]{ref-fey2019fast}
Fey, Matthias, e Jan Eric Lenssen. 2019. {``Fast Graph Representation
Learning with PyTorch Geometric''}. \emph{ICLR Workshop on
Representation Learning on Graphs and Manifolds}.

\bibitem[\citeproctext]{ref-forrester2008engineering}
Forrester, Alexander, Andras Sobester, e Andy Keane. 2008.
\emph{Engineering Design via Surrogate Modelling: A Practical Guide}.
Wiley.

\bibitem[\citeproctext]{ref-gilmer2017neural}
Gilmer, Justin, Samuel S. Schoenholz, Patrick F. Riley, Oriol Vinyals, e
George E. Dahl. 2017. {``Neural Message Passing for Quantum
Chemistry''}. \emph{International Conference on Machine Learning},
1263--72.

\bibitem[\citeproctext]{ref-goodfellow2016deep}
Goodfellow, Ian, Yoshua Bengio, e Aaron Courville. 2016. \emph{Deep
Learning}. MIT Press.

\bibitem[\citeproctext]{ref-goodfellow2014generative}
Goodfellow, Ian, Jean Pouget-Abadie, Mehdi Mirza, et al. 2014.
{``Generative Adversarial Nets''}. \emph{Advances in Neural Information
Processing Systems}.

\bibitem[\citeproctext]{ref-gray2019openmdao}
Gray, Justin S., John T. Hwang, Joaquim R. R. A. Martins, Kenneth T.
Moore, e Bret A. Naylor. 2019. {``OpenMDAO: An open-source framework for
multidisciplinary design, analysis, and optimization''}.
\emph{Structural and Multidisciplinary Optimization} 59: 1075--104.

\bibitem[\citeproctext]{ref-ho2020denoising}
Ho, Jonathan, Ajay Jain, e Pieter Abbeel. 2020. {``Denoising Diffusion
Probabilistic Models''}. \emph{Advances in Neural Information Processing
Systems}.

\bibitem[\citeproctext]{ref-jasa2018open}
Jasa, John P., John T. Hwang, e Joaquim R. R. A. Martins. 2018.
{``Open-source coupled aerostructural optimization using Python''}.
\emph{Structural and Multidisciplinary Optimization} 57: 1815--27.

\bibitem[\citeproctext]{ref-karniadakis2021physics}
Karniadakis, George Em, Ioannis G. Kevrekidis, Lu Lu, Paris Perdikaris,
Sifan Wang, e Liu Yang. 2021. {``Physics-informed machine learning''}.
\emph{Nature Reviews Physics} 3: 422--40.

\bibitem[\citeproctext]{ref-kingma2013auto}
Kingma, Diederik P., e Max Welling. 2013. \emph{Auto-Encoding
Variational Bayes}. \url{https://arxiv.org/abs/1312.6114}.

\bibitem[\citeproctext]{ref-kipf2017semi}
Kipf, Thomas N., e Max Welling. 2017. {``Semi-Supervised Classification
with Graph Convolutional Networks''}. \emph{International Conference on
Learning Representations}.

\bibitem[\citeproctext]{ref-lecun2015deep}
LeCun, Yann, Yoshua Bengio, e Geoffrey Hinton. 2015. {``Deep
learning''}. \emph{Nature} 521: 436--44.

\bibitem[\citeproctext]{ref-martins2021engineering}
Martins, Joaquim R. R. A., e Andrew Ning. 2021. \emph{Engineering Design
Optimization}. Cambridge University Press.

\bibitem[\citeproctext]{ref-raissi2019physics}
Raissi, Maziar, Paris Perdikaris, e George E. Karniadakis. 2019.
{``Physics-informed neural networks: A deep learning framework for
solving forward and inverse problems involving nonlinear partial
differential equations''}. \emph{Journal of Computational Physics} 378:
686--707.

\bibitem[\citeproctext]{ref-rawlings2017model}
Rawlings, James B., David Q. Mayne, e Moritz Diehl. 2017. \emph{Model
Predictive Control: Theory, Computation, and Design}. Nob Hill
Publishing.

\bibitem[\citeproctext]{ref-rombach2022high}
Rombach, Robin, Andreas Blattmann, Dominik Lorenz, Patrick Esser, e
Bjorn Ommer. 2022. {``High-Resolution Image Synthesis with Latent
Diffusion Models''}. \emph{IEEE/CVF Conference on Computer Vision and
Pattern Recognition}, 10684--95.

\bibitem[\citeproctext]{ref-zhou2020graph}
Zhou, Jie, Ganqu Cui, Shengding Hu, et al. 2020. {``Graph neural
networks: A review of methods and applications''}. \emph{AI Open} 1:
57--81.

\end{CSLReferences}




\end{document}
